\section{Discussion and Limitations}

\subsection{Applicability Conditions}
\label{sec:discussion-drift}

SAKI applies when co-located RocksDB instances share a fixed background budget,
write pressure drifts across tenants at window time scales, and the host can
observe tenant-local demand and maintenance symptoms. The repeated epoch runs
establish the main placement effect, the continuous runs show that it survives
runtime rate-limiter actuation, and the mechanism checks rule out stale skew,
tenant-local autotuning, pressure-only tracking, lower-layer throttling, and a
non-maintenance actuation path. Ranking-margin replay and perturbation runs
further show that the main continuous high-budget assignment is fixed by a
demand-aware score family rather than by one isolated coefficient vector.
Together, these results support the systems claim that a host-level
rank-and-assign controller can place an existing compaction/I/O budget better
than equal sharing.

The bounded-drift condition matters. SAKI relies on previous-window
observations remaining relevant to the next window. The main continuous
workload moves the high-pressure set by one tenant per 20-second window, and
\fastdrift pushes this to 1.5 tenants/window while retaining
CI-strict high-pressure benefit with overlap 2.57/4. Abrupt reversals,
sub-window bursts, and within-window phase changes may outrun the controller.
The 20-second window, 4/8/4 tier split, and concrete budget levels are
mechanism-isolation settings. In the continuous anchor, a 10/20/40-second
window sweep keeps the 96\,MB/s budget fixed and brackets the trade-off:
10 seconds improves assignment frequency but makes high-P99 and LOW collateral
noisier, whereas 40 seconds reduces overlap and leaves adaptation unbounded in
all three seeds. The conservation rule in Section~\ref{sec:saki-design}
supports other tier counts, budget spreads, or a continuous allocator.

\subsection{Deployment Assumptions}

The continuous harness changes background rate limits through \ratecall. The
measured rate-limiter coverage microstudy shows this knob shapes
flush/compaction I/O in our configuration, so foreground latency and
throughput changes arise indirectly from maintenance service rather than
direct foreground throttling. RocksDB's built-in auto-tuned rate limiter
remains useful local machinery, but it adapts within per-tenant ceilings;
SAKI supplies the host-level loop that moves unused budget across tenants. A
natural deployment variant is to let SAKI choose tenant ceilings while RocksDB
autotuning adapts inside each ceiling; we did not evaluate that combined
policy.

A production deployment would need tenant inventory, periodic RocksDB
statistics, a conserved host budget, and per-tenant runtime actuation. The
offered-demand signal should be measured near admission, such as at a client
proxy, RPC ingress path, tenant queue, or write-admission counter, before
upstream backpressure, drops, or retries hide true pressure. Operators would
likely add smoothing, hysteresis, cooldowns, minimum dwell times, and SLO-aware
floors around the simple rank-and-assign rule. If metrics or tenant inventory are
missing, the safe operational choice is to retain the last valid assignment or
fall back to all-mid sharing, preserving the aggregate budget before
re-enabling reallocation. The reference score should be treated as a
demand-aware instance of the framework, not as a portable constant vector:
deployments with different admission signals or maintenance dynamics should
renormalize terms while preserving budget conservation and demand observability.
Other LSM engines need analogous tenant-scoped maintenance signals and a
runtime maintenance-budget knob.

\subsection{Resource and Fairness Boundaries}

SAKI controls one resource dimension: background I/O budget placement. CPU,
cache, memory, device queueing, and per-level write amplification may become
the limiting resources in other deployments. A deployment can detect this
regime when rate-limiter traffic stops tracking flush/compaction output or
when CPU, cache, or device-queue metrics saturate while compaction debt remains
high. The continuous efficiency metric is therefore normalized bytes/write
because SAKI completes more foreground writes; a full maintenance-cost
decomposition is a separate multi-resource control problem.

The policy is also a targeted reallocation. Tier-wise reporting makes the
collateral cost explicit because shifting budget to high-pressure tenants can
cost lower-pressure tenants. In the main continuous aggregate, LOW tenants see
small write-tail and throughput costs while HIGH tenants improve substantially;
the 4\,KB-value workload leaves P999 not CI-separated from zero despite a
CI-strict P99 win.
These results motivate SLO-aware floors and fairness policy as deployment
layers around the same fixed-budget loop. The read-heavy variant raises
point-read QPS around the main harness; scan-heavy workloads, range-scan SLOs,
and read-dominant deployments require their own scoring and validation layer.

\subsection{Trace Validation Boundary}

The public-trace audits define the conditions required for meaningful
controller validation. CacheLib provides real cache-front-end request traces
and validates trace-shaped replay plumbing, but the selected segments create
too little RocksDB maintenance stress: completed logical writes were only
about 5--6\% of offered demand and compaction output plateaued near 30\,MB/s.
Baleen exposes a different issue: aggregate scaled write supply can look
adequate while too few tenants are active in each window to actuate the fixed
all-mid budget. Across 146 independent candidates, no segment had enough
static-actuatable supply to justify a second static smoke.

The next production-trace validation layer needs storage-engine traces with
tenant identity, drift, sustained write pressure, and actuatable supply under
the same fixed-budget contract. Non-leveled compaction styles, longer
deployments, and joint multi-resource control are natural extensions after
that trace condition is met.
